\chapter{Estymacja i predykcja bayesowska: od posterioru do decyzji}
\label{ch:bayes-estimation-prediction}
\paragraph{Powiązania i odsyłacze.} Konstrukcje posteriorów pochodzą z rozdziałów~\ref{ch:bayes-foundations} i~\ref{ch:conjugate-families}; funkcje straty z rozdziału~\ref{ch:decision}.


\section{Co właściwie chcemy oszacować?}
W rozdziałach~\ref{ch:bayes-foundations}--\ref{ch:conjugate-families}
wyprowadziliśmy rozkłady a posteriori. Sam posterior nie rozstrzyga
jednak, \emph{jaką liczbę powinniśmy podać} ani \emph{jaką decyzję
podjąć}. Te dwa pytania mają różne odpowiedzi: pierwsze wymaga
wyboru funkcji straty, drugie dodatkowo określenia rzeczywistych
konsekwencji działań. Jeszcze inne pytanie dotyczy prognozy
nowej obserwacji. Dla przykładu w kontroli jakości rozróżniamy
trzy cele: oszacowanie nieznanego odsetka wad $p$, przewidzenie
liczby wad w następnej partii i podjęcie decyzji o zatrzymaniu produkcji.

\begin{definicja}[Reguła bayesowska i ryzyko posteriorowe]
\label{def:bayes-rule-ch19}
Niech $(\Theta_0,\mathcal T)$, $(\mathsf X,\mathcal X)$ i
$(\mathsf A,\mathcal G)$ będą przestrzeniami mierzalnymi,
$\pi$ właściwym priorem, $K(\theta,\cdot)$ jądrem modelu danych,
a $\Pi(x,\cdot)$ jego posteriorem. Niech
$L:\Theta_0\times\mathsf A\to[0,\infty]$ będzie mierzalną
funkcją straty. Dla obserwacji $x$ i działania $a$ definiujemy
\emph{ryzyko posteriorowe}
$\rho_x(a)=\int_{\Theta_0}L(\theta,a)\,\dd\Pi(x,\theta)$.
Mierzalną regułę $\delta:\mathsf X\to\mathsf A$ nazywamy
\emph{regułą Bayesa}, jeśli $\rho_x(\delta(x))=\inf_{a\in\mathsf A}\rho_x(a)$
dla $M$-prawie każdego $x$, gdzie $M$ jest rozkładem
predykcyjnym danych przed ich obserwacją. Istnienie mierzalnego
minimizatora wymaga osobnego uzasadnienia; sama definicja go
nie gwarantuje.
\end{definicja}

\begin{twierdzenie}[Optymalizacja posteriorowa minimalizuje ryzyko Bayesa]
\label{thm:bayes-risk-tower}
W założeniach definicji~\ref{def:bayes-rule-ch19} niech
$\delta_*$ będzie mierzalną regułą Bayesa. Wtedy dla każdej
mierzalnej reguły $\delta:\mathsf X\to\mathsf A$ zachodzi
\[
\int_{\Theta_0}\int_{\mathsf X}
 L(\theta,\delta_*(x))\,\dd K(\theta,x)\,\dd\pi(\theta)
\le
\int_{\Theta_0}\int_{\mathsf X}
 L(\theta,\delta(x))\,\dd K(\theta,x)\,\dd\pi(\theta).
\]
Dla obu stron dopuszczamy wartość $+\infty$.
\end{twierdzenie}
\begin{proof}
Nieujemność i mierzalność straty pozwalają zastosować Tonellego.
Z własności warunkowania, najpierw dla funkcji prostych, a potem
przez monotoniczną zbieżność,
\[
\int_{\Theta_0}\int_{\mathsf X}L(\theta,\delta(x))
\,\dd K(\theta,x)\,\dd\pi(\theta)
=\int_{\mathsf X}\rho_x(\delta(x))\,\dd M(x).
\]
Przy każdym $x$ poza zbiorem $M$-miary zero reguła $\delta_*$
minimalizuje wyrażenie pod całką. Całkowanie tej nierówności
daje tezę. To minimalizacja ryzyka \emph{uśrednionego względem
prioru}; nie wynika z niej dominacja ryzyka częstotliwościowego
przy każdym ustalonym $\theta$. Wynik dotyczy wyłącznie
porównywanych reguł, ustalonego modelu i przyjętej funkcji straty.
\end{proof}

\section{Funkcja straty wyznacza estymator}
W twierdzeniu~\ref{thm:bayes-actions} wykazaliśmy, że strata
kwadratowa prowadzi do średniej posterioru, a bezwzględna do
mediany. Do decyzji o rezerwie lub zapasie części zamiennych
bardziej pasuje niekiedy \emph{strata asymetryczna}: zaniżenie
przewidywanej liczby awarii kosztuje więcej niż jej zawyżenie.

\begin{definicja}[Asymetryczna strata liniowa]
Dla $\tau\in(0,1)$ i rzeczywistych $\theta,a$ definiujemy
$L_\tau(\theta,a)=\tau(\theta-a)_++(1-\tau)(a-\theta)_+$,
gdzie $t_+=\max(t,0)$. Parametr $\tau$ określa względny
koszt niedoszacowania, a $1-\tau$ przeszacowania.
\end{definicja}

\begin{twierdzenie}[Kwantyl posterioru jako optymalna decyzja]
\label{thm:bayes-quantile-loss}
Niech $\Pi_x$ będzie rozkładem na $\R$ o skończonym
pierwszym momencie, a $F_x(t)=\Pi_x(({-}\infty,t])$ jego
dystrybuantą. Dla $\tau\in(0,1)$ liczba $a$ minimalizuje
$\int_\R L_\tau(\theta,a)\,\dd\Pi_x(\theta)$ po $a\in\R$
wtedy i tylko wtedy, gdy
$F_x(a-)\le\tau\le F_x(a)$.
W szczególności przy ciągłej, ściśle rosnącej dystrybuancie
jedyną optymalną decyzją jest kwantyl $F_x^{-1}(\tau)$.
\end{twierdzenie}
\begin{proof}
Strata jest wypukła względem $a$ i ograniczona przez
$|\theta|+|a|$, więc jej wartość oczekiwana jest skończona.
Dla $a$ niebędącego atomem rozkładu pochodna po $a$ jest
równa $F_x(a)-\tau$: przy $\theta<a$ pochodna wynosi
$1-\tau$, a przy $\theta>a$ jest równa $-\tau$.
Ogólnie jednostronne pochodne funkcji ryzyka to
$F_x(a-)-\tau$ oraz $F_x(a)-\tau$, co wynika
z ilorazów różnicowych i zbieżności zdominowanej.
Wypukła funkcja osiąga minimum w punkcie $a$ dokładnie
wtedy, gdy lewa pochodna jest nie większa od zera,
a prawa nie mniejsza od zera. To daje tezę.
\end{proof}

\begin{przyklad}[Wielkość zapasu części na kolejny tydzień]
Niech predykcyjna liczba awarii $Y$ ma rozkład, dla którego
kwantyl rzędu $0{,}9$ wynosi $q_{0{,}9}$. Jeśli koszt
brakującej części jest dziewięciokrotnie większy niż koszt
zbędnej części, strata liniowa ma $\tau=0{,}9$.
Optymalna liczba części jest \emph{całkowitoliczbowym}
kwantylem predykcyjnym rzędu $0{,}9$, a nie średnią
posterioru intensywności. W praktyce należy również sprawdzić
koszty magazynowania, ograniczenia pojemności i możliwość
dokupienia części w trakcie tygodnia.
\end{przyklad}

\section{Rozkład parametru a rozkład przyszłych obserwacji}
Niech $\lambda$ oznacza intensywność awarii na godzinę,
a nowe zliczenie $Y_*$ dotyczy $e_*>0$ godzin. Gdy znamy
$\lambda$, losowość $Y_*$ jest poissonowska. Gdy jej nie
znamy, predykcja musi scałkować tę losowość po posteriorze.

\begin{twierdzenie}[Predykcja gamma--Poisson ma większą wariancję]
\label{thm:gamma-poisson-dispersion}
Jeżeli $\lambda\mid X=x\sim\operatorname{Gamma}(r,v)$,
$r,v>0$, a $Y_*\mid\lambda,X=x\sim\operatorname{Pois}(e\lambda)$
dla $e>0$, to
\[
\E[Y_*\mid X=x]=\frac{er}{v},\qquad
\Var(Y_*\mid X=x)=\frac{er}{v}+\frac{e^2r}{v^2}.
\]
W szczególności $\Var(Y_*\mid X=x)>\E[Y_*\mid X=x]$.
\end{twierdzenie}
\begin{proof}
Korzystamy kolejno z
$\E[Y_*\mid\lambda,x]=e\lambda$,
$\Var(Y_*\mid\lambda,x)=e\lambda$ i z prawa całkowitej
wariancji względem posterioru. Pierwszy składnik daje
$e\E[\lambda\mid x]=er/v$, drugi
$e^2\Var(\lambda\mid x)=e^2r/v^2$.
Dodatniość ostatniego składnika wynika z $e,r,v>0$.
\end{proof}

\begin{przyklad}[Ile awarii przewidzimy?]
Z rozdziału~\ref{ch:conjugate-families} weźmy posterior
$\operatorname{Gamma}(6,300)$ (intensywność na godzinę).
Przy ekspozycji $e=50$ predykcja ma średnią $1$ i wariancję
$7/6$. Dokładny rozkład liczby awarii jest ujemny
dwumianowy; z twierdzenia~\ref{thm:gamma-poisson}
$\Prob(Y_*=0\mid x)=(6/7)^6\approx0{,}39657$,
a $\Prob(Y_*\ge3\mid x)\approx0{,}09356$.
Jeśli konieczne jest zabezpieczenie przed co najmniej trzema
awariami, stosujemy ostatnie prawdopodobieństwo, a nie tylko
średnią prognozy równą jeden.
\end{przyklad}

\section{Predykcja optymalna względem właściwych funkcji oceny}
Przy ocenianiu prognoz probabilistycznych chcemy nagradzać
podawanie całego rozkładu, a nie tylko liczby. Oczekiwane
ryzyko obliczamy zawsze względem jawnie wskazanego prawa
przyszłej obserwacji; nie jest to automatycznie to samo, co
częstościowa ocena modelu na niezależnej próbie testowej. \emph{Właściwa
funkcja oceny} (ang. \emph{proper scoring rule}) jest taką funkcją
straty, dla której prawdziwy rozkład predykcyjny minimalizuje
oczekiwaną stratę. Rozróżnijmy dwa elementarne przypadki.

\begin{twierdzenie}[Poprawność miary Briera dla zdarzenia binarnego]
\label{thm:brier}
Niech $Y\in\{0,1\}$ przy zadanej informacji ma rozkład
$\operatorname{Bernoulli}(p)$, $p\in[0,1]$. Prognozujemy
$p$ liczbą $q\in[0,1]$ i stosujemy stratę Briera
$S(q,Y)=(q-Y)^2$. Wtedy
\[
\E[S(q,Y)]=(q-p)^2+p(1-p).
\]
Jedynym optymalnym wyborem jest $q=p$.
\end{twierdzenie}
\begin{proof}
Ponieważ $Y^2=Y$, mamy
$\E(q-Y)^2=q^2-2qp+p=(q-p)^2+p-p^2$.
Zatem nadwyżka ryzyka względem prawidłowej prognozy
jest dokładnie kwadratem błędu prognozowanego
prawdopodobieństwa.
\end{proof}

\begin{twierdzenie}[Reguła logarytmiczna i dywergencja KL]
\label{thm:log-score-bayes}
Niech $\mathsf Y=\{1,\ldots,k\}$, $k\ge2$, a
$P=(p_1,\ldots,p_k)$ będzie rzeczywistym rozkładem przyszłej
obserwacji. Dla prognozy $Q=(q_1,\ldots,q_k)$ z $q_j>0$
gdy $p_j>0$, strata logarytmiczna ma postać
$S(Q,Y)=-\log q_Y$. Wówczas
\[
\E_P S(Q,Y)-\E_P S(P,Y)
=\sum_{j:p_j>0}p_j\log\frac{p_j}{q_j}
= D_{\mathrm{KL}}(P\Vert Q)\ge0,
\]
a równość zachodzi dokładnie, gdy $Q=P$.
\end{twierdzenie}
\begin{proof}
Pierwsza równość to bezpośrednie odjęcie skończonych sum.
Do dowodu nierówności stosujemy $\log t\le t-1$ do
$t=q_j/p_j$ dla $p_j>0$:
$\sum_{p_j>0}p_j\log(q_j/p_j)
\le\sum_{p_j>0}(q_j-p_j)\le0$.
Jeśli zachodzi równość, to każdy składnik wymusza
$q_j=p_j$ dla $p_j>0$; normalizacja wymusza też
$q_j=0$ poza nośnikiem $P$.
\end{proof}

\begin{uwaga}[Walidacja predykcji]
Porównując modele na wspólnym zbiorze testowym, możemy
obliczyć średnią stratę Briera dla prognoz prawdopodobieństw
lub średnią stratę logarytmiczną. Nie należy dobierać priorów
i parametrów na tym samym zbiorze, który ma następnie służyć
za niezależną ocenę uogólnienia. Przy szeregach czasowych
konieczne jest zachowanie chronologii walidacji.
\end{uwaga}

\section{Przedział wiarygodności, przedział ufności i przedział predykcyjny}
\begin{przyklad}[Trzy różne pytania dla siedmiu wad na sto]
Załóżmy $S=7$, $n=100$ w modelu Bernoulliego. Estymator MLE
odsetka wad to $\widehat p=0{,}07$. Klasyczny 95\% przedział
Wilsona z rozdziału~\ref{ch:confidence-intervals} jest w przybliżeniu
$[0{,}03432,0{,}13750]$. Przy priorze
$\operatorname{Beta}(1,1)$ posterior to
$\operatorname{Beta}(8,94)$, którego centralny 95\%
przedział wiarygodności wynosi w przybliżeniu
$[0{,}03482,0{,}13760]$.
Oba przedziały są liczbami bliskie, ale oznaczają co innego:
pierwszy ma (przy przyjętym modelu) przybliżoną gwarancję
powtarzalnego pokrycia, drugi ma posteriorową masę $0{,}95$
po obserwacji danych. Predykcja liczby wad w kolejnych
dziesięciu sztukach nie jest żadnym z tych przedziałów:
ma rozkład beta--dwumianowy, a
$\Prob(T=0\mid S=7)=B(8,104)/B(8,94)\approx0{,}45772$.
Do planowania liczby kontroli lub zapasu części należy używać
rozkładu predykcyjnego, a nie wyłącznie niepewności parametru.
\end{przyklad}

\section{Sprawdzanie modelu przez replikację danych}
Analiza posterioru zakłada, że określony model generowania
danych jest wystarczająco wiarygodny. \emph{Posterior predictive
check} (PPC) polega na wylosowaniu $\theta\sim\Pi(\cdot\mid x)$,
a potem nowych danych $X^{\rm rep}\sim P_\theta$ oraz
porównaniu statystyki diagnostycznej $T(X^{\rm rep})$ z $T(x)$.
Statystyka $T$ powinna celować w charakterystyczne odchylenie:
liczbę zer, rozrzut, ogony, autokorelację lub zależność między
grupami. Nie jest to automatycznie klasyczny test o jednostajnym
rozkładzie wartości p.

\begin{przyklad}[Dlaczego posterior predictive p-value nie jest jednostajna?]
Niech $p\sim\operatorname{Beta}(1,1)$ i
$X\mid p\sim\operatorname{Bernoulli}(p)$, zaś
$X^{\rm rep}\mid p$ niech będzie warunkowo niezależną kopią.
Jako statystykę przyjmijmy $T(X)=X$ i zdefiniujmy
$q(X)=\Prob(X^{\rm rep}\ge X\mid X)$.
Dla $X=0$ zachodzi $q(0)=1$. Dla $X=1$
posterior jest $\operatorname{Beta}(2,1)$, więc
$q(1)=\E[p\mid X=1]=2/3$.
W rozkładzie predykcyjnym przed obserwacją
$\Prob(X=0)=\Prob(X=1)=1/2$, stąd $q(X)$ przyjmuje
wyłącznie wartości $1$ i $2/3$, każdą z prawdopodobieństwem
$1/2$. Nie jest jednostajna na $[0,1]$.
To konkretny powód, aby nie interpretować automatycznie PPC
jako standardowej wartości p z gwarancją poziomu testu.
\end{przyklad}

\begin{stwierdzenie}[Rozkład predykcyjny beta--dwumianowy przed obserwacją]
\label{prop:prior-predictive-betabinomial}
Niech $p\sim\operatorname{Beta}(a,b)$, $a,b>0$,
a warunkowo przy znanym $p$ zmienna $S$ ma rozkład
$\Bin(n,p)$, $n\ge1$. Wtedy dla $s=0,\ldots,n$
\[
\Prob(S=s)=\binom ns\frac{B(a+s,b+n-s)}{B(a,b)}.
\]
Przy priorze jednostajnym $\operatorname{Beta}(1,1)$
rozkład $S$ jest jednostajny na $\{0,\ldots,n\}$.
\end{stwierdzenie}
\begin{proof}
Pierwszy wzór wynika z całkowania
$\binom ns p^s(1-p)^{n-s}$ względem gęstości prioru,
co daje iloraz całek beta. Dla $a=b=1$ korzystamy
z $B(s+1,n-s+1)=s!(n-s)!/(n+1)!$ i otrzymujemy
$\binom ns B(s+1,n-s+1)=1/(n+1)$.
\end{proof}
\begin{uwaga}
Jednostajność liczby sukcesów przed obserwacją \emph{nie oznacza},
że wszystkie wartości $p$ są równie wiarygodne po zobaczeniu danych.
Rozkład predykcyjny przed obserwacją może ujawnić, czy prior
przewiduje realistyczne dane: w modelu beta--dwumianowym
$\operatorname{Beta}(1,1)$ dopuszcza bardzo duży rozrzut liczby
wad, mimo że pojedyncza obserwacja ma prawdopodobieństwo
wady $1/2$. Przy aktualnej wiedzy o rzadkich awariach
taki prior może być nieodpowiedni.
\end{uwaga}

\section{Praktyczny schemat: od celu do metody}
\begin{center}\small
\begin{tabular}{>{\raggedright\arraybackslash}p{.27\textwidth}>{\raggedright\arraybackslash}p{.31\textwidth}>{\raggedright\arraybackslash}p{.27\textwidth}}
\toprule
\textbf{Pytanie} & \textbf{Właściwe narzędzie} & \textbf{Co koniecznie sprawdzić}\tabularnewline
\midrule
Jaka jest nieznana wartość parametru? & Średnia lub kwantyl posterioru, zależnie od straty & Rozkład a priori i identyfikowalność.\\
Ile awarii będzie w kolejnej partii? & Rozkład posterior predictive, jego ogon i kwantyle & Ekspozycja i zmienność między okresami.\\
Jak dobrać zapas? & Kwantyl predykcyjny przy asymetrycznej stracie & Rzeczywiste koszty braków i nadmiaru.\\
Który model lepiej prognozuje? & Brier lub log-loss na danych testowych & Przeciek danych i zależność czasowa.\\
Czy model jest wiarygodny? & PPC, diagnostyka reszt, walidacja predykcyjna & PPC nie jest automatycznie testem poziomu $\alpha$.\\
Czy mamy gwarancję pokrycia dla każdego parametru? & Klasyczny przedział ufności, np. Wilsona & Rodzaj gwarancji i założenia losowania.\\
\bottomrule
\end{tabular}
\end{center}

\Needspace{26\baselineskip}
\begin{lstlisting}[language=Python,caption={Kompletny program: estymacja, predykcja, ryzyko alarmu i przedziały (Python, SciPy).}]
from scipy.stats import beta, betabinom, nbinom

n, defects = 100, 7
a, b = 1, 1
post_a, post_b = a + defects, b + n - defects
print("Posterior mean:", post_a / (post_a + post_b))
print("95% credible interval:", beta.ppf([0.025, 0.975], post_a, post_b))

# Defects in the next ten independent devices.
future_n = 10
print("P(no defects):", betabinom.pmf(0, future_n, post_a, post_b))
print("P(at least two defects):", betabinom.sf(1, future_n, post_a, post_b))
print("Predictive 90% quantile:", betabinom.ppf(0.9, future_n, post_a, post_b))

# Failures in 50 hours, given posterior Gamma(shape=6, rate=300).
shape, rate, exposure = 6, 300, 50
prob_3plus = nbinom.sf(2, shape, rate / (rate + exposure))
print("P(three or more failures):", prob_3plus)

# A warning costs 3 if unjustified; a missed event costs 50.
false_alarm, missed_event = 3.0, 50.0
threshold = false_alarm / (false_alarm + missed_event)
print("Alarm threshold:", threshold, "Warn:", prob_3plus > threshold)
\end{lstlisting}

\begin{zadanie}[Rezerwa dla nierównej straty]
W posteriorze $p\mid x\sim\operatorname{Beta}(8,94)$ chcemy
podać jedną liczbę reprezentującą $p$. Za jednostkową wartość
niedoszacowania płacimy dziewięć razy więcej niż za przeszacowanie.
Jaki estymator Bayesa stosujemy?
\par\emph{Rozwiązanie.} Korzystamy z asymetrycznej straty
liniowej o $\tau=9/(9+1)=0{,}9$. Twierdzenie
\ref{thm:bayes-quantile-loss} daje $F_{\operatorname{Beta}(8,94)}^{-1}(0{,}9)$.
Średnia $8/102$ minimalizowałaby inną, kwadratową stratę,
więc nie jest tu rozwiązaniem zadania decyzyjnego.
\end{zadanie}

% ================================================================
% ROZDZIAL 20 -- BAYESOWSKIE TESTOWANIE HIPOTEZ
% ================================================================
